\qquad Em várias situações do cotidiano é necessário manter conjuntos de
dados organizados: da lista telefônica ao histórico de transações
bancárias, a ordenação é uma operação central em ciência da
computação. Diante disso, diversos algoritmos de ordenação foram
desenvolvidos ao longo dos anos, cada um com características
próprias de desempenho, simplicidade de implementação e
comportamento em diferentes cenários de entrada.

\qquad Este trabalho é dividido em cinco capítulos, um para cada
algoritmo de ordenação estudado: o Capítulo 1 trata do Insertion
Sort, o Capítulo 2 do Selection Sort, o Capítulo 3 do Bubble Sort, o
Capítulo 4 do Shell Sort e o Capítulo 5 do Merge Sort. Os quatro
primeiros são algoritmos de ordenação por comparação amplamente
utilizados como introdução ao estudo de algoritmos e possuem
complexidade $O(n^2)$ no pior caso (no caso do Shell Sort, com a
sequência de intervalos original utilizada neste trabalho, e apenas
para entradas construídas especificamente para isso) --- ainda que,
dependendo dos dados de entrada, alguns deles possam ser bem mais
rápidos do que essa cota sugere. O Merge Sort, por sua vez, baseia-se
no paradigma de dividir para conquistar e tem complexidade
$\Theta(n \log n)$ em todos os casos, ao custo de utilizar memória
adicional proporcional ao tamanho da entrada.

\qquad Ao longo de cada capítulo são apresentados o funcionamento do
respectivo algoritmo, um teste de mesa manual, a análise de
complexidade nas diferentes situações de entrada, e os resultados
experimentais obtidos a partir de uma implementação em C++, testada
com instâncias de tamanhos crescentes (de 10 a 1.000.000 de
elementos) e três padrões de entrada: crescente, decrescente e
randômica. Ao final do trabalho, os cinco algoritmos são comparados
entre si por meio de uma tabela e um gráfico geral, seguidos de uma
conclusão consolidada sobre os resultados observados.
